\documentclass[aps,pra,twocolumn,superscriptaddress,longbibliography]{revtex4-2}

\usepackage{amsmath,amssymb,amsthm,graphicx,xcolor,bm,braket,dcolumn,booktabs}
\usepackage{hyperref}

\newtheorem{theorem}{Theorem}
\newtheorem{corollary}{Corollary}
\newtheorem{prediction}{Prediction}
\newtheorem{observation}{Observation}

\newcommand{\sthree}{S^{3}}

\begin{document}

\title{Angular Sparsity Is an Atomic-Sector Property:\\
The Abelian Residue of Multi-Center Bonding}
\thanks{Paper 58 in the Geometric Vacuum series}

\author{J.~Loutey}
\affiliation{Independent Researcher, Kent, Washington}

\date{August 10, 2026}

%% ========================================================================
%% ABSTRACT
%% ========================================================================

\begin{abstract}
The \textsc{GeoVac} program builds molecular Hamiltonians from
atom-centered hydrogenic blocks (screened $Z_{\mathrm{eff}}$ per block,
\emph{not} shared-$p_{0}$ Coulomb--Sturmian, which Papers~8--9 rule out
for heteronuclear pairs) whose angular matrix elements are
generated from quantum-number labels by Gaunt selection, with no spatial
quadrature.  This yields structurally sparse qubit Hamiltonians---838
Pauli terms for LiH at $Q = 30$, exactly linear in $Q$ across
molecules at fixed basis (Paper~14)---at
the cost of a documented binding failure for second-row diatomics
(the W1e wall, Papers~19/20~\cite{loutey_paper19,loutey_paper20}).

We settle the question of whether that trade was necessary.  Using exact
rational two-center machinery (Mulliken auxiliary functions $A_n$/$B_n$
for overlaps; a hydrogenic eigen-trick for cross-center one-body terms,
certified by an exact bra/ket identity holding entrywise), we census the
\emph{genuine} bare tensor $(S, h, g)$ for a real fragment pair and find:
\textbf{[MEASURED]} the Gaunt $l$-selection rule does not survive
cross-center evaluation, while the axial $m$-rule does; cross-center
one-body entries reach $0.80$~Ha---bond scale, not perturbative.
\textbf{[COUNTED]} the symmetry-\emph{permitted} electron-repulsion density
inflates $13.8\times$ over the builder at $n_{\max} = 2$ and $15.0\times$
at $n_{\max} = 3$; the $n_{\max} = 2$ figure is independently corroborated
numerically ($29.8\%$ measured against $29.4\%$ counted), the
$n_{\max} = 3$ figure is counted only.

We give the mechanism \textbf{[INTERNAL THEOREM]}: $m$ labels the
\emph{abelian} symmetry shared by both centers (rotation about the
internuclear axis is one global one-parameter group, so $m$ is a common
label), while $l$ labels irreducible representations of a rotation group
attached to a single center.  The intersection of the two centers'
rotation groups is exactly the axial $U(1)$, whose character group is
$\mathbb{Z}$ by Pontryagin duality.  Abelian selection rules are additive
charge conservation and survive; non-abelian Clebsch--Gordan selection
requires the full sphere and does not.  The surviving \emph{cross-center}
angular rule is therefore the one generated by the common axial subgroup,
which is readable from the molecular point group in advance.

Two re-readings follow.  First, the $17.9\times$ Pauli inflation
previously attributed to L\"owdin orthogonalization
(v4.73.1) was not, in the bulk, orthogonalization's price: the genuine
tensor's own permitted density inflates by the same order ($\sim 15\times$)
before any transformation, so L\"owdin was paying a bill that was largely
already owed.  Second \textbf{[OBSERVATION]}, four previously attempted
relocations of the non-orthogonality---into the operator, into a dual
basis, into a high-symmetry embedding, into device-measured many-body
overlaps---each moved the cost without reducing it, because the cost
measures the physical asymmetry of a two-center problem.  A fifth,
measured here, is the sharpest: enabling equivalent-atom-swap tapering
\emph{reduces} the qubit saving ($\mathrm{BeH_2}$ $12 \to 9$,
$\mathrm{N_2}$ $22 \to 19$, $\mathrm{F_2}$ $22 \to 15$) while inflating
the Pauli count $2.3$--$3.6\times$, because the swap rotation merges
per-block abelian gradings into joint stabilizers.  A permutation
identifies blocks, and identifying blocks destroys the per-block labels
that generate the sparsity.

We close with a demonstration \textbf{[PANEL-VERIFIED]}: all-electron
12-electron NaH, the framework's hardest documented failure, binds at
$R_{\rm eq} = 3.736\,a_0$ ($+4.8\%$ vs experiment) with $91.1\%$ of
in-basis FCI binding from three fragment-native determinants, once the
integrals are genuine end-to-end---with \emph{fewer} correlation degrees of
freedom than the runs that failed, which is what makes the inference go
through.  This localizes W1e to the Hamiltonian specification rather than
the correlation treatment on this system.  The demonstration deliberately
forfeits the sparsity: it makes no sparsity claim.

\textbf{[SYMBOLIC + MEASURED]} Separately, the diatomic integrals are made
\emph{quadrature-free}:\ the minimal-basis H$_2$ FCI energy is certified to
$60$ digits with symbolic $\tau$-series termination, its entire potential
curve $E(R)$ is a single closed-form expression in $\{\exp, E_{1}, \log, \gamma\}$
with certified $R_{\rm eq}$, $D_{e}$ and $k$, and the heteronuclear LiH case
reaches $30$ digits net of an explicit tail bound.

A by-product bears on the continuous theory the census borrows from:\ the
radial integral at the core of the two-center direct-Coulomb $(AA|BB)$
repulsion closes in elementary form (weight one, $\pi$-free, no
dilogarithm), so that block is not merely counted but \emph{decided}---all $195$
permitted entries proven nonzero by Lindemann--Weierstrass separation.
\end{abstract}

\maketitle

%% ========================================================================
\section{Introduction: what the sparsity bought, and what it cost}
\label{sec:intro}
%% ========================================================================

The \textsc{GeoVac} framework's computational proposition is structural
sparsity.  Angular matrix elements are generated from the quantum-number
labels $(n, l, m)$ by Gaunt selection rules and Wigner $3j$ symbols with
no spatial integration, so the resulting second-quantized Hamiltonian is
sparse by construction rather than by numerical thresholding.  The
downstream consequence is the qubit-resource result of
Paper~14~\cite{loutey_paper14}: composed LiH at 838 Pauli terms on
$Q = 30$ qubits, a within-molecule basis exponent of $O(Q^{3.17})$
against $O(Q^{3.9\text{--}4.3})$ for Gaussian baselines --- a narrower
margin than the retired pair-diagonal figures suggested --- and a
$76\times$ reduction against cc-pVDZ.

The framework's documented failure is chemical binding for second-row
diatomics.  Named W1e, it has resisted an unusually thorough sequence of
attacks: six Phillips--Kleinman modifications, basis-level Schmidt
orthogonalization, core-correlation estimates, basis enlargement to
$n_{\max} = 4$, an explicit-core Hartree--Fock treatment with twelve
electrons, and a DMRG cross-check that reproduced FCI bit-identically at
finite sector dimension without producing an interior minimum
(Papers~19/20~\cite{loutey_paper19,loutey_paper20}; CHANGELOG v2.19.4--v3.87.0).  The wall was localized to
the Hamiltonian specification---the projection from the continuous
Level-4 solver to second-quantized $(h_1, \text{eri}, e_{\rm core})$
integrals---rather than to the treatment of correlation.

This paper answers the question that localization leaves open.  The
framework neglects the overlap between orbitals on different centers,
treating a non-orthogonal basis as though it were orthonormal, and that
neglect is what preserves the sparsity.  Was there an alternative?  Could
one keep the genuine cross-center physics \emph{and} the label-generated
sparsity, by declining to transform the integral tensor at all?

The answer is no, and the reason is group-theoretic rather than technical.

\subsection{A change of method, not a refinement}
\label{sec:method_break}

\emph{Relation to the continuous theory.}  Evaluating integrals over
genuine two-center orbital products is, in itself, nothing new:\ it is
standard Coulomb--Sturmian practice, developed over decades by Avery and
co-workers on the many-center hyperspherical-harmonic
foundation~\cite{avery1989, avery2006}, and the framework's cross-center
one-body couplings are precisely the Shibuya--Wulfman
integrals~\cite{shibuya1965} of that theory
(Paper~19~\cite{loutey_paper19}).  \textsc{GeoVac} adds nothing to that
continuous apparatus.  What is new is not the integrals but the
\emph{discrete} reading of them:\ which of the continuous theory's
couplings survive as exact, label-forced zeros when the angular content
is carried to a qubit Hamiltonian, and which do not.  In that sense this
paper is the discrete/combinatorial skeleton of a continuous
subject---the relationship the angular selection rules already bear to
their Gaunt and Wigner origins, now traced into the two-center case.

It is worth stating plainly what this paper does differently, because the
step is easy to mistake for an incremental improvement in integral quality.
It is not.  Every previous \textsc{GeoVac} molecular builder evaluates its
integrals in the \emph{atomic sector}: each orbital product entering an
integral sits on a single nucleus, and inter-nuclear coupling enters
through a surrogate.  Three generations, distinguished by what they do at
the bond:

\begin{enumerate}
\item \textbf{Composed geometry} (Paper~17~\cite{loutey_paper17}).  A
  Phillips--Kleinman pseudopotential replaces explicit two-body
  core--valence Coulomb with one-body $Z_{\mathrm{eff}}$ screening, so
  \emph{cross-block electron-repulsion integrals vanish identically}.  The
  bond is not present in the two-body tensor at all.

\item \textbf{Balanced coupled composition}
  (Papers~19/20~\cite{loutey_paper19,loutey_paper20}).  The pseudopotential
  is eliminated and cross-block repulsion is restored, balanced by a
  cross-center nuclear-attraction term.  But both restored pieces remain
  atomic-sector objects: the cross-center one-body integral is
  $\langle \psi^{A}_{nlm} | -Z_B/|\mathbf{r}-\mathbf{R}_B| |
  \psi^{A}_{n'l'm'} \rangle$, with \emph{both} orbitals on center $A$ and
  only the nucleus displaced, and the restored cross-block repulsion is
  constructed so as to continue obeying Gaunt triangle selection.  The
  orbital products are still one-center; the $l$-selection rule still holds.

\item \textbf{This paper.}  The integrals are evaluated over orbital
  products that genuinely span two nuclei, $\langle \chi^{A} | \cdot |
  \chi^{B} \rangle$.  That is the step neither prior generation takes.
\end{enumerate}

The distinction is not cosmetic, and the census below measures its price
exactly.  Because generations 1 and 2 never form a two-center orbital
product, the Gaunt $l$-selection rule is exact for them by construction,
and the label-generated sparsity follows.  Once the product genuinely
spans two centers, that rule fails---only the axial $m$-rule survives---and
the permitted density inflates by $13.8\times$ at $n_{\max} = 2$.  The
sparsity was never a discovered property of molecular electronic structure;
it was a property of declining to evaluate the bond.

The demonstration in Sec.~\ref{sec:nah} is the corresponding gain.  NaH is
the system the prior generations could not bind under any modification
attempted, and it binds here with three fragment-native determinants.  We
therefore read W1e not as a defect to be repaired within the atomic-sector
construction, but as the accumulated cost of that construction---which
makes the present route a replacement for it rather than a correction to
it.  What the replacement forfeits is precisely the qubit-resource
proposition Paper~14 rests on; Sec.~\ref{sec:rereadings} states that
trade as a conservation law rather than a defeat.

\emph{Scope of this subsection.}  The three-generation description is a
statement about what the builders compute, verifiable from their
construction, not a new result.  The demonstration below (Sec.~\ref{sec:nah})
is a diatomic.  A polyatomic extension has since been measured
(BeH$_2$ and H$_2$O, basis ladder to $M = 19$): at matched minimal basis and
matched correlation treatment the present route improves the geometry only
modestly over generation~1 ($+10.0\%$ vs $11.7\%$ for BeH$_2$, $+14.9\%$ vs
$19.4\%$ for H$_2$O), but the error is basis-limited rather than
method-limited and falls to $+1.4\%$ and $+2.5\%$ respectively once the basis
is enlarged---and the H$_2$O bond angle, an observable generation~1 cannot
represent at all, comes out at $106.3^\circ$ against an experimental
$104.5^\circ$.  Those figures are RHF over Gaussian-fitted Slater shapes in
the same engine used for Sec.~\ref{sec:nah}; they are reported here as scope,
not as a result of this paper.

%% ========================================================================
\section{The exact-arithmetic census}
\label{sec:census}
%% ========================================================================

\subsection{Machinery and certificate}

Two-center overlaps at arbitrary $(n, l, m, Z)$ are computed as exact
rationals via the Mulliken auxiliary functions
$A_n$/$B_n$~\cite{mulliken1949,roothaan1951}---one of several classical
routes to two-center Slater integrals~\cite{barnett_coulson1951,
filter_steinborn1978}---evaluated in prolate spheroidal coordinates
where the two-center problem
separates~\cite{loutey_paper11}.  Because the result is rational,
vanishing is \emph{decidable} rather than inferred from a numerical
tolerance---a property no quadrature scheme can offer, and the reason
this census can make a claim about exact zeros at all.

Cross-center one-body elements are obtained by extending the same route
with $1/r_A$ and $1/r_B$ kernels, each of which cancels against the
volume element $(\xi^2 - \eta^2)$, followed by the hydrogenic eigen-trick.
The assembly carries its own certificate: the bra/ket identity
\begin{equation}
E_a S - Z_B I_B \;=\; E_b S - Z_A I_A
\label{eq:braket_certificate}
\end{equation}
holds as an \emph{exact rational equality} on a six-case panel and on every
entry of the census configuration---and coefficient-wise on both
exponential legs, which is strictly stronger than agreement at a single
$R$.  Its diagnostic power is itself tested rather than assumed:\ two
independent corruptions of the inputs (a wrong bra eigenvalue, and $I_A$
interchanged with $I_B$) each break the identity, so it is demonstrably not
an algebraic tautology.  Four independent
quadrature cross-checks agree at $\sim 2 \times 10^{-16}$.

\subsection{Result}

The census configuration is $Z_A = 3$ at the origin, $Z_B = 1$ at
$R = 3\,\hat{z}$, hydrogenic orbitals to $n_{\max} = 2$ on each center,
$M = 10$ spatial orbitals.  Table~\ref{tab:census} compares the genuine
bare tensor against the production builder's.

\begin{table}[b]
\caption{Genuine versus builder sparsity, $Z_A=3$/$Z_B=1$ at
$R = 3\,a_0$, $n_{\max} = 2$ per center ($M = 10$).  ``Genuine zeros
are'' records which selection rule accounts for the surviving zeros.}
\label{tab:census}
\begin{ruledtabular}
\begin{tabular}{lrrrl}
object & genuine & builder & dense & genuine zeros are \\
\hline
$S$ & 32 & 10 & 100 & $m$-rule $+$ same-center \\
$h$ & 44 & 22 & 100 & $m$-rule only \\
$g$ & 2{,}944 & 214 & 10{,}000 & $m$-rule $+$ same-center Gaunt \\
\end{tabular}
\end{ruledtabular}
\end{table}

\emph{Provenance, stated precisely because the three rows differ, and
because the rows are internally mixed.} The cross-center blocks of $S$ and
$h$ (22 entries each) are evaluated in exact rational arithmetic by the
machinery of Sec.~\ref{sec:census}\,A, so their \emph{non}-vanishing is
decided rather than inferred.  The same-center $S$ block is fixed
analytically by orthonormality of same-$Z$ hydrogenics rather than
recomputed; the same-center $h$ block is evaluated with the production
float routine, so its zeros there are threshold decisions.  And the
vanishing of every non-matching-$m$ entry is not evaluated at all---it is
forced by Theorem~\ref{thm:abelian_residue}(i).  In this census the
rational machinery decides non-vanishing; the zeros come from symmetry.  The $g$ row is \emph{counted}: it is a
symmetry-rule census in the complex-$m$ basis, because no two-center
electron-repulsion engine exists in the framework---the $m$-rule and
per-side Gaunt feasibility are applied combinatorially, and no two-center
electron-repulsion integral was numerically evaluated.  The $g$ figures
are therefore claims about how many entries the selection rules
\emph{permit}, not verifications that each permitted entry is nonzero.
Supplying that engine is the natural follow-on to this paper.

Every count in Table~\ref{tab:census} is re-derived from the selection
rules alone in \texttt{tests/test\_paper58\_census.py}.  Both columns
apply the \emph{same} axial rule $m_p + m_r = m_q + m_s$; the builder
column is the genuine column restricted to quartets whose four
orbitals share one center, with the single-center shared-$L$ Gaunt
condition imposed identically in both (the same $176$ removals at
$n_{\max}=2$).  The inflation factor therefore measures the
two-center span and nothing else---it is not an artifact of
comparing a pair-diagonal $m$-convention against a global one,
since neither column uses the pair-diagonal rule.

\emph{Status of that follow-on.}  It is now partially supplied, and the
per-class outcome bears directly on what the $g$ row could become.  Three
of the four classes are closed in exact form and validated against
independent engines: the one-center classes, $(AA|BB)$, and the hybrid
pair $(AA|AB)$/$(AB|BB)$.  For the exchange class $(AB|AB)$---Ruedenberg's
own subject~\cite{ruedenberg1951}---the assembly is validated for general
$(l,m)$ and its ordered radial double integral, the last non-closed step,
now closes as well, for general $(\tau, \sigma, H, j)$ and unequal
exponents (agreement with nested quadrature at $10^{-13}$).  All four
classes therefore have closed-form machinery.

\begin{center}
\begin{tabular}{lll}
\hline\hline
class & closed form & transcendental content \\
\hline
one-center, $(AA|BB)$ & yes & none---elementary \\
$(AA|AB)$, $(AB|BB)$  & yes & $\{E_1, \ln\}$ \\
$(AB|AB)$             & yes & $\{E_1, \ln, \gamma\}$ \\
\hline\hline
\end{tabular}
\end{center}

The transcendental content grows across the classes, but---and this is the
structurally important part---it does \emph{not} escalate in weight.  The
ordered radial integral is an iterated integral over a simplex, the shape
that defines a period, and iterated integrals of weight-one objects
generically land at weight two (dilogarithms, $\zeta(2)$).  This one does
not: assembled in closed form it agrees with quadrature to
$6.3 \times 10^{-16}$ and contains only $\exp$, $E_1$, $\ln$ and Euler's
$\gamma$.  Nor does $\sigma \neq 0$ break it.  The derivatives
$d^{\sigma} Q_{\tau}/d\xi^{\sigma}$ carry poles of order up to $\sigma$ at
$\xi = \pm 1$, but they are absorbed exactly by the $(\xi^2-1)^{H}$
prefactor, because per electron the net exponent there is
$(|m_a| + |m_b| - |m_a - m_b|)/2 \geq 0$ by the triangle inequality---with
equality, i.e.\ exact absorption, whenever $m_a$ and $m_b$ have opposite
signs.  Verified to $2.5 \times 10^{-14}$ with no weight-two content at
any $\sigma$ tested.  So the seed sets above form a weight-one filtration
throughout.

\textbf{This does not upgrade the $g$ row}, for two independent reasons.
First, no re-census has been run with the new machinery, so the counted
figures stand exactly as counted.  Second, and structurally, what the row
needs is zero-\emph{decidability}, and that is available only where a
closed form exists.

For $(AA|BB)$ that argument now goes through, and by exactly the route the
$S$ and $h$ rows already use.  Every $\pi$ cancels from the closed form---
the harmonic normalisations against the $4\pi/(2L+1)$ of the multipole
potential---leaving

\begin{equation}
  (ab|cd) \;=\; A_0(R) \;+\; \sum_j A_j(R)\, e^{-\lambda_j R},
  \label{eq:aabb_closed_form}
\end{equation}

with $A_j$ rational functions and $\lambda_j$ rational.  Checked on
$(1s\,1s|1s\,1s)$, $(2p_0 2p_0|1s\,1s)$ and $(2p_{+1} 2p_{+1}|2p_0 2p_0)$
across $Z$ combinations: the $\pi$ power is common (indeed zero) in every
term.  For algebraic $R$ the exponents $-\lambda_j R$ are distinct
algebraic numbers, so by Lindemann--Weierstrass the family
$\{1, e^{-\lambda_j R}\}$ is linearly independent over the algebraic
numbers and the expression vanishes \emph{iff} every $A_j$ does.  That is
the same separation the $S$ and $h$ rows are decided by, applied to a
class that now has the closed form to support it.

The hybrid and exchange classes do not follow: they carry
$\{E_1, \ln\}$ and $\{E_1, \ln, \gamma\}$ respectively, and would need an
independence argument over those seeds, which we have not attempted.

\emph{Consequence for this paper's own sector.}  The classes carrying
residual $l$-selection---$(AA|AA)$, $(BB|BB)$, $(AA|BB)$---are precisely
the ones that are now zero-\emph{decidable}: the one-centre classes
trivially---they are exact rationals---and $(AA|BB)$ by the argument
above.  So the sector that Theorem~\ref{thm:abelian_residue} makes a
positive prediction about is no longer merely counted but decidable.

\emph{And the decision has now been run on $(AA|BB)$.}  On the census
configuration, that block holds $625$ entries of which the counting rules
permit $195$.  Deciding each one by the separation above---grouping the
closed form by exponential rate and asking whether every rational
coefficient vanishes---gives

\begin{center}
\begin{tabular}{lr}
\hline\hline
outcome & count \\
\hline
decided nonzero & 195 \\
zero, vanishing Gaunt coefficient & 0 \\
zero, accidental (Gaunt nonzero, radial sum cancels) & 0 \\
\hline\hline
\end{tabular}
\end{center}

\textbf{All $195$ permitted entries are genuinely nonzero.}  There are no
accidental zeros and no symmetry zeros that the counting missed, so on
this block the counted density is not an overcount but is \emph{exactly}
the true density.  This is a decision rather than a measurement: no
threshold is involved anywhere, and it upgrades the block from
\textsc{counted} to \textsc{decided}.  The same decision run over a
$5\times 7$ grid of charge pairs and internuclear distances---including the
equal-charge and rate-coincidence configurations, where degenerate
exponential rates are the likeliest source of an accidental cancellation,
and a node-bearing $n_{\max}=3$ sample---returns $195/195$ nonzero at every
point, so the verdict is not an artifact of the census
configuration.

It also settles, for this block, the reading the Gaussian corroboration
could only suggest.  That check found $29.8\%$ against $29.4\%$ counted
and attributed the gap to the real-Cartesian versus complex-$m$ basis
convention rather than to over-permissive counting; the decided census
confirms there is no over-permissiveness here to find.

The $g$ row's overall tier is nonetheless unchanged, because $(AA|BB)$ and
the one-centre classes together are only the $20.5\%$ that carries
residual $l$-selection.  The three genuinely cross classes---$79.5\%$ of the
tensor---remain counted, and will until an independence argument over
$\{E_1, \ln, \gamma\}$ is available.

We state the outlook carefully, because both natural readings are wrong.
The pessimistic one---that the seed set escalates---is wrong: it is a
weight-one filtration at every class, and a closed form now exists for every
class.  But the optimistic one---that the remaining independence argument is
therefore \emph{ordinary work}---is wrong too, because \emph{weight one is
not the same as decidable}.  The $(AA|BB)$ decision went through only because
that class is the one place $\pi$ cancels \emph{and} no $E_1$, $\ln$ or
$\gamma$ appears: pure $\{e^{-\lambda R}\}$, which Lindemann--Weierstrass
separates at any algebraic $R$.  The cross classes carry $E_1(\mu R)$ (hybrid
and exchange) and, for the exchange class, a standalone Euler $\gamma$ that
survives assembly at every $\tau$.  Deciding an entry's vanishing at a fixed
$R$ then requires those seed \emph{values} to be linearly independent---and
$E_1$-value independence at algebraic points is an open problem of
transcendence theory, while $\gamma$ is not known even to be irrational.
That is a genuine obstruction in principle, not ordinary work, and it draws
a \emph{categorical} line---$(AA|BB)$ decidable, the $E_1/\gamma$ classes
not---rather than the gradient the weight-one filtration alone suggests.

What \emph{is} available now is one rung weaker and unconditional.  The seed
\emph{functions} $\{1, e^{-\lambda R}, E_1(\mu R), \ln R\}$ are linearly
independent over $\mathbb{Q}(R)$---\textbf{[SYMBOLIC (proof-by-argument)]} a
theorem (Liouville--Rosenlicht: $E_1$ is
the standard non-elementary integral, transcendental over the exponential
field)---so every cross-class entry is nonzero \emph{as a function of $R$}
unless its grouped coefficients vanish, and any zero is therefore isolated at
special $R$.  With the Gaussian corroboration above already showing the
permitted entries generically nonzero, the honest position for the cross
classes is: functionally nonvanishing (unconditional), generically nonzero
(measured), and pointwise-\textsc{decidable} only \emph{conditional} on a
Schanuel-type independence hypothesis for the $E_1$ and $\gamma$ values.  That
is materially better than counted, and blocked from decided by open
transcendence rather than by missing labour.

\emph{Beyond two centres (scope).}  The transcendence does not merely climb the
weight ladder at the third centre; it changes ladder.  The two-body
three-centre integral $(XY|XZ)$---two two-centre densities sharing one centre
on different axes, the first object a polyatomic requires and one no prolate
spheroidal system can hold (that coordinate has exactly two foci)---is
naturally evaluated in momentum space, where the Coulomb kernel is $4\pi/k^2$
and each centre is a phase rather than a focus.  There the two densities
Fock-project onto spheres of \emph{different} radii $\sqrt{c_1}\neq\sqrt{c_2}$,
and the product of their stereographic weights is the quartic
$y^2=(c_1 k^2+1)(c_2 k^2+1)$---an elliptic curve.  Its $D=0$ period is exactly a
complete elliptic integral, and the full integral is a genus-one object on the
same family---by Paper~59's final characterization an irreducible $\Gamma(2)$
period (the elementary elliptic-polylogarithm target is excluded there)---%
degenerating to the genus-zero two-centre case only
on the measure-zero locus $c_1=c_2$ of a shared sphere.  The two negatives that
make the two-centre engine weight-one and $(AA|BB)$ decidable---no $\pi$, no
dilogarithm---are therefore \emph{two-centre} properties: the third centre
leaves the polylogarithm hierarchy for genus-one elliptic transcendentals.
This is the transcendence-level form of the reason Gaussian orbitals, whose
product theorem collapses three centres to one, displaced Slater orbitals for
polyatomic quantum chemistry.  It is reported here as scope; the method and the
elliptic identification are numerical (momentum-space evaluation validated
against an independent Gaussian reference to $\sim\!10^{-14}$, and the elliptic
period verified to $31$ digits), not a closed form.

What the build does settle is narrower, and worth stating because it is
exactly the sector this paper's mechanism concerns.  The classes in which
the residual $l$-selection survives---$(AA|AA)$, $(BB|BB)$ and
$(AA|BB)$, together $20.5\%$ of the genuine tensor---are now exactly
computable in closed form and carry no transcendental seed whatever.  An
$l$-selection census restricted to those classes is therefore no longer
gated on the missing engine.  That is the sharpest thing the build
supplies for this paper: the sector where
Theorem~\ref{thm:abelian_residue} predicts residual $l$-selection to
survive is now exactly computable, and computable with no transcendental
seed at all, so a census over it would test the theorem's positive content
directly rather than by counting.

\emph{The two halves compose.}  The census machinery of
Sec.~\ref{sec:census} supplies exact two-centre $S$ and $h$---the
Mulliken/Ruedenberg auxiliary route and the hydrogenic eigen-trick---and
named the missing two-centre electron-repulsion engine as its one gap.
With that engine supplied, the three ingredients now compose into a
complete calculation carrying no Gaussian fit anywhere.  Verified on
$\mathrm{H}_2$ at $R = 1.4$~bohr in a $1s$-per-centre basis: total FCI
energy $-1.106556606091$~Ha.  Evaluating the identical basis through
Gaussian-fitted integrals converges monotonically onto that value as the
fit deepens ($2.7\times10^{-3}$, $3.0\times10^{-4}$, $4.5\times10^{-5}$,
$8.0\times10^{-6}$, $1.9\times10^{-6}$~Ha at $4, 6, 8, 10, 12$
primitives), confirming the native value is the exact one.  The residual
sits in $h$, not in $g$: the electron-repulsion tensors agree to
$1.9\times10^{-8}$ and the overlaps to $1.5\times10^{-8}$ at ten
primitives, while the one-electron terms lag because kinetic energy probes
the nuclear cusp, where a sum of Gaussians converges slowly by
construction.

Subject to that distinction, the electron-repulsion tensor is $29.4\%$
permitted-dense genuinely against $2.1\%$ for the builder.

\emph{Numerical corroboration of the $g$ row.} The counting can be checked
without the missing Neumann engine~\cite{ruedenberg1951}, by evaluating
the same tensor with a
McMurchie--Davidson Gaussian engine~\cite{mcmurchie1978} over fitted
Slater shapes.  On the
census configuration this gives $2{,}976$ of $10{,}000$ entries above
$10^{-10}$, i.e.\ $29.8\%$ dense against the $29.4\%$ counted---agreement
at the level the basis convention permits, since the numerical evaluation
is in the real Cartesian basis and the count is in the complex-$m$ basis,
which differ by the $\pm m$ mixing.  So the permitted entries are
\emph{generically nonzero}: the density is a real property of the tensor
and not an artifact of over-permissive counting.  The
$l$-selection failure appears at bond scale in individual entries---for
instance $(1s_A\,2p_{z,B} \,|\, 2s_A\,2s_A) = -0.218$, where the
same-center analog vanishes identically.

This does not upgrade the row to \textsc{measured}: Gaussian-fitted floats
make these threshold decisions rather than decided zeros, so the exact
statement still awaits the engine.  What it removes is the possibility
that the counted density was an overcount.

\emph{The $m$-rule, verified without reference to $m$.}
Theorem~\ref{thm:abelian_residue}(i) has a basis-free consequence that can
be tested directly on the tensor: since $1/r_{12}$ is rotationally
invariant and both nuclei lie on the axis, a rotation about the
internuclear axis maps the basis into itself and must leave the tensor
fixed.  It does, to $2.2 \times 10^{-16}$.  Displacing the nuclei off that
axis and repeating the identical check gives $1.3 \times 10^{-1}$---fifteen
orders of magnitude larger.  The symmetry is axial, and only axial.  At $n_{\max} = 3$ ($M = 28$) the comparison is
$114{,}280$ \emph{permitted} entries ($18.6\%$) against $7{,}600$
($1.2\%$): the inflation factor is $13.8\times$ at $n_{\max} = 2$ and
$15.0\times$ at $n_{\max} = 3$, i.e.\ \emph{increasing} with basis size.
Both are counted; only the $n_{\max} = 2$ figure is numerically
corroborated below.

Two further measurements matter for interpretation.  First, per-side
Gaunt feasibility removes $5.6\%$ of entries at $n_{\max} = 2$ and
$6.3\%$ at $n_{\max} = 3$, beyond what the $m$-rule already removes, and the three cross classes
$(AA|AB) + (AB|AB) + (AB|BB)$ account for $79.5\%$ of the genuine tensor at $n_{\max}=2$ ($80.0\%$ at $n_{\max}=3$).
The residual $l$-selection is therefore confined to the classes carrying a
same-center distribution on each side---$(AA|AA)$, $(BB|BB)$ and
$(AA|BB)$---together $20.5\%$ of the genuine tensor; the three genuinely
cross classes carry none.  Second, cross-center magnitudes reach bond
scale:\ the largest overlap is $0.52$ and the largest one-body entry is
$0.80$~Ha.  These are \emph{maxima}, not a range---the census block also
contains much smaller entries (the smallest overlap is $0.011$, on the
$(2p_{0}, 2p_{0})$ pair).  An earlier statement of them as intervals
$0.08$--$0.52$ and $0.19$--$0.80$ took its lower ends from an inconsistent
subset of the census pairs and was incorrect.  For
scale, the entire NaH bond is under $2$~eV, so individual entries in the
dense part of the tensor exceed the binding energy being computed.  They
are not droppable as small corrections---an independent quantification of
the W1d/W1e diagnosis.

%% ========================================================================
\section{Mechanism: the abelian residue}
\label{sec:mechanism}
%% ========================================================================

The census result is not an accident of the chosen configuration.  It is
forced, and the forcing is elementary once stated in the right language.

\begin{theorem}[Angular selection survives molecularization only in the
abelian sector]
\label{thm:abelian_residue}
Let two nuclei $A$, $B$ lie on the $z$-axis, with atom-centered basis
functions labeled $(n, l, m)$ about their own center, and let
$\hat{O}$ be invariant under the molecular symmetry group.  Then:
\begin{enumerate}
\item[(i)] The $m$-selection rule holds exactly across centers:
$\langle ab | \hat{O} | cd \rangle = 0$ unless
$m_a + m_b = m_c + m_d$.
\item[(ii)] No $l$-selection rule holds between basis functions on
different centers.  More strongly, the cross-center support is \emph{full}
in $l$ at fixed $m$:\ the re-expansion does not terminate, so every
admissible $l$ couples.
\item[(iii)] Of the per-center $SO(3)$ labels, only those of the common
axial subgroup $U(1) \cong SO(2)$ survive cross-center; by Pontryagin
duality its characters are the integers, so the surviving \emph{cross-center}
rule is additive $m$-conservation.
\end{enumerate}
\end{theorem}

Part (iii) is deliberately narrow, and two limits are worth stating because
the wider reading is false.  It does not say the axial $U(1)$ generates all
surviving sparsity: same-center Gaunt selection survives from the full
per-center $SO(3)$ (visible in Table~\ref{tab:census}'s $g$ and $S$ rows),
and the largest group fixing both nuclei is not $U(1)$ but $C_{\infty v}$,
which is \emph{non-abelian}---its $\sigma_v$ reflection is exactly the
Hopf $m_l \to -m_l$ $\mathbb{Z}_2$ used as the baseline of
Table~\ref{tab:swap}.  So a non-abelian common group does yield usable
label structure.  What (iii) claims is only that the \emph{continuous}
per-center angular labels reduce to the axial one across centers, which is
what the census measures.

\begin{proof}[Proof sketch]
For (i): rotation about the internuclear axis is a \emph{single}
one-parameter group acting on all of $\mathbb{R}^3$.  Its generator
$L_z$ is independent of which point on the axis is taken as origin, so
$m$ is a global label shared by both centers rather than a per-center
one.  Invariance of $\hat{O}$ under this group gives additive
conservation of $m$.

For (ii): $l$ labels irreducible representations of the rotation group
about a \emph{specific} center.  Rotations about $A$ and rotations about
$B$ are distinct subgroups of the Euclidean group whose intersection is
precisely the axial subgroup.  A function that is an $l$-eigenfunction
about $B$, re-expanded about $A$, has support on all $l$: the two-center
re-expansion does not terminate.  Hence no $l$-selection between cross-center
pairs.  (The same-center blocks retain Gaunt selection, consistent with
the $5.6\%$--$6.3\%$ residual measured in Sec.~\ref{sec:census}.)

For (iii): the intersection in (ii) is $SO(2) \cong U(1)$, abelian.
Abelian compact groups have one-dimensional irreducibles labeled by
characters, and by Pontryagin duality the character group of $U(1)$ is
$\mathbb{Z}$---which is the integer label $m$.  Selection rules from
abelian symmetry are additive conservation laws on integer charges, and
are preserved by any operation respecting the axis.  Non-abelian
selection rules are Clebsch--Gordan couplings requiring the full group;
reducing the group destroys them.
\end{proof}

The theorem is independently corroborated by the framework's own
production code, along a path that knows nothing of this analysis: the
cross-center nuclear-attraction routine in
\texttt{shibuya\_wulfman} returns
zero by an explicit early exit when $m_1 \neq m_2$, and imposes no
$l$-selection on the element: its per-multipole Gaunt conditions (triangle
inequality and $l_1 + L + l_2$ parity) are satisfiable by some $L$ for
every $(l_1, l_2)$, so no $l$ pair is excluded.  The $m$-rule is hard-coded because it
holds; the $l$-rule is absent because it does not.

Theorem~\ref{thm:abelian_residue} converts a measurement into a
prediction, because the largest common abelian symmetry is readable from
the molecular point group before any integral is computed.

\begin{corollary}[Sparsity budget from the point group]
\label{cor:point_group}
For a molecule with point group $G$, label-generated angular sparsity is
bounded by the grading available from the abelian structure of $G$.
Linear molecules ($C_{\infty v}$, $D_{\infty h}$) carry a continuous
$U(1)$ and therefore a $\mathbb{Z}$ grading, worth a multiplicative
factor.  Bent molecules carry only a finite point group and therefore a
finite grading.
\end{corollary}

$\mathrm{H_2O}$ has point group $C_{2v}$, which is order four and
\emph{abelian}, isomorphic to $\mathbb{Z}_2 \times \mathbb{Z}_2$; all its
irreducibles are one-dimensional, so Corollary~\ref{cor:point_group}
predicts a two-bit spatial grading---two qubits, not a multiplicative
factor.  We attempted to test this and could not, for a reason that is
itself a structural finding.

\begin{observation}[The composed builder carries no angular geometry]
\label{obs:no_angle}
\texttt{molecular\_spec.py} contains no bond angle: \texttt{h2o\_spec} is
\texttt{hydride\_spec(8)}, parameterized by nuclear charge and a
\emph{single} O--H distance distributed across five blocks (core, two
bond pairs, two lone pairs), with \texttt{MolecularSpec.nuclei = None}.
The composed $\mathrm{H_2O}$ Hamiltonian therefore does not represent
$C_{2v}$ at all, and the framework's reported $\mathrm{H_2O}$ geometry
error is a purely radial coordinate with no angular degree of freedom.
Corollary~\ref{cor:point_group} is consequently \emph{not well posed} on
this builder in the bent case; supplying hand-built bent nuclei would
test a symmetry the Hamiltonian does not encode.
\end{observation}

\begin{prediction}[Conditional, for any angular build]
\label{pred:angular}
In a build that does carry angular geometry, a molecule whose point group
is abelian of order $2^k$ admits a $k$-bit spatial grading and no more.
This remains falsifiable, but not on the present builder.
\end{prediction}

What \emph{is} measurable on the present builder is the linear case, where
geometry is angle-free and therefore unambiguous---and it returns a result
we did not anticipate, reported as
Observation~\ref{obs:permutation_cost} below.

The measured factor of $3$--$5$ for the $m$-rule at the bases of
Sec.~\ref{sec:census} is the $\mathbb{Z}$-grading case of
Corollary~\ref{cor:point_group}.  It is exact and structural, and it is a
constant factor---not the exactly-linear-in-$Q$ Pauli-count story
($N_{\mathrm{Pauli}} = 27.90 \times Q$ across molecules at fixed
basis, Paper~14) the framework's qubit results rest on.

%% ========================================================================
\section{Two re-readings}
\label{sec:rereadings}
%% ========================================================================

\subsection{L\"owdin was paying an honest bill}

The corpus previously recorded a $17.9\times$ Pauli inflation from
retrofitting L\"owdin orthogonalization~\cite{lowdin1950} onto the
balanced LiH integrals
(CHANGELOG v4.73.1), and a $14\times$ inflation in the earlier
heterogeneous-nested experiment.  Both were read as the price of
orthogonalization, and orthogonalization was accordingly treated as the
sparsity-destroying move.

The census shows this reading was wrong in an instructive way.  The
genuine tensor's own permitted-density inflation over the builder is
$\sim 15\times$---the same order as the $17.9\times$ L\"owdin figure,
though the two are not commensurable: $17.9\times$ is a Pauli-term count on
L\"owdin-retrofitted balanced LiH integrals, $15.0\times$ a counted
permitted-entry ratio on a $Z_A = 3$/$Z_B = 1$ hydrogenic census, and
neither carries an error bar.  The inference does not need them to agree
numerically: the \emph{bulk} of that inflation cannot have been
orthogonalization's price, because the untransformed tensor never possessed
the sparsity L\"owdin was said to destroy.  ``L\"owdin destroys sparsity''
and ``the molecular tensor is intrinsically $l$-dense cross-center'' are
the same fact observed twice, from two sides.

This matters because it removes an entire class of proposed remedies.
Any scheme whose premise is ``avoid the transformation and keep the
sparsity'' is answering a question that was mis-posed.

\subsection{Cost conservation}

\begin{observation}[Relocating the metric does not reduce it]
\label{obs:cost_conservation}
Five independent attempts to move the non-orthogonality out of the way
each succeeded in moving it and failed to reduce it:
\begin{enumerate}
\item into the \textbf{operator} (L\"owdin $S^{-1/2}$): operator
densifies $17.9\times$;
\item into a \textbf{dual basis} (biorthogonal encoding): densifies
comparably.  The metric provably does not disappear---in the contravariant
representation it reappears as an $S^{-1}$ contraction, and the dual
creation/annihilation operators cease to be Hermitian
conjugates~\cite{artacho1991}---but whether that $S^{-1}$ is
\emph{dense} is a property of the two-center Slater basis, measured here
rather than imported;
\item into a \textbf{high-symmetry embedding} (hyperspherical Level~4N):
angular basis requirement explodes---LiH at $63.5\%$ $R_{\rm eq}$ error
with $l_{\max} = 2$ genuinely insufficient;
\item into \textbf{device-measured many-body overlaps} (NOCI): succeeds
at the state level, but the operator is dense independently;
\item into a \textbf{permutation symmetry} (equivalent-atom swap): net
loss, measured---see Observation~\ref{obs:permutation_cost}.
\end{enumerate}
The invariant is that the cost measures the physical asymmetry of a
two-center problem.  The currency is a choice; the total is not.
\end{observation}

The fifth entry was measured for this paper and is the sharpest instance,
because the trade is visible term by term.

\begin{observation}[Permutation symmetry costs qubits on this builder]
\label{obs:permutation_cost}
Enabling equivalent-atom-swap tapering \emph{reduces} the qubit saving and
inflates the Pauli count, in every swap-eligible configuration tested
(Table~\ref{tab:swap}).  The mechanism is visible in the retained
stabilizer labels: the swap rotation merges the per-block Hopf and
$\ell$-parity stabilizers into joint \texttt{SWAP\_JOINT} stabilizers, so
one permutation stabilizer is bought at the price of several per-block
abelian gradings.  \textbf{[MEASURED]}
\end{observation}

\begin{table}[b]
\caption{Equivalent-atom-swap tapering against two baselines.
$\Delta Q$ is qubits removed (larger is better); Pauli counts are for the
tapered operator.  LiH has no equivalent atoms and is the null control.}
\label{tab:swap}
\begin{ruledtabular}
\begin{tabular}{llrrr}
system & baseline & $\Delta Q$ & $\Delta Q$ +swap & Pauli ratio \\
\hline
$\mathrm{BeH_2}$ & Hopf only & 7 & 6 & $3.50\times$ \\
$\mathrm{BeH_2}$ & Hopf $+\ \ell$ & 12 & 9 & $3.60\times$ \\
$\mathrm{N_2}$ & Hopf $+\ \ell$ & 22 & 19 & $2.30\times$ \\
$\mathrm{F_2}$ & Hopf only & 12 & 9 & $3.50\times$ \\
$\mathrm{F_2}$ & Hopf $+\ \ell$ & 22 & 15 & $3.60\times$ \\
LiH & Hopf $+\ \ell$ & 8 & 8 & $1.00\times$ \\
\end{tabular}
\end{ruledtabular}
\end{table}

This is consistent with Theorem~\ref{thm:abelian_residue} rather than
incidental to it.  A permutation is not an abelian grading of the existing
block structure; it is an operation that \emph{identifies} blocks, and
identifying blocks destroys the per-block labels that generated the
sparsity.  The abelian residue is not merely what survives---it is what
can be exploited without collapsing the structure that carries it.

\emph{Documentation discrepancy, since corrected.} The docstring of
\texttt{extended\_tapered\_from\_spec} previously stated that atom-swap ``saves
additional qubits on polyatomics ($\mathrm{BeH_2}$ +1, $\mathrm{H_2O}$ +1,
$\mathrm{NH_3}$ +1, $\mathrm{CH_4}$ +2).''  The measured $\mathrm{BeH_2}$
value is $-1$ against the Hopf-only baseline and $-3$ against Hopf
$+\ \ell$; the sign is wrong under both.  We report the measurement and
do not attempt to adjudicate whether the documentation was always
incorrect or the behavior regressed.

Observation~\ref{obs:cost_conservation} is labeled an observation, not a
theorem: it is a pattern across five instances with a plausible cause,
not a proof that no sixth currency exists.  The underlying obstruction
has been characterized in the framework as non-commuting center
projections, with $\|[P_A, P_B]\| = 0.50$ at LiH equilibrium and
principal angles $7.6^\circ$/$44.7^\circ$/$67.3^\circ$.

%% ========================================================================
\section{Demonstration: NaH binds with genuine integrals}
\label{sec:nah}
%% ========================================================================

If the wall is the Hamiltonian specification, then supplying genuine
end-to-end integrals should bind NaH \emph{without needing more}
correlation than the failing runs used.  It does---with considerably less.

Two scope conditions govern what follows, and the argument is weaker
without them.  First, this construction \textbf{makes no sparsity claim}:
the radial functions are six-Gaussian fits of Slater shapes evaluated with
a McMurchie--Davidson engine~\cite{mcmurchie1978} and a dense
non-orthogonal Slater--Condon solver, so it forfeits precisely the
label-generated sparsity this paper is about---it is item~4 of
Observation~\ref{obs:cost_conservation}, not an escape from it.  Second,
the correlation treatment is not held fixed: the framework's failing runs
were FCI and DMRG in an orthogonal second-quantized space, whereas this is
three non-orthogonal determinants.  That asymmetry is what makes the
inference go through---binding appears with \emph{less} correlation, so
correlation cannot have been what was missing.

The construction is all-electron 12-electron NaH with
$\mathrm{Na}\{1s, 2s, 2p \times 3, 3s\}$ and $\mathrm{H}\{1s\}$
($M = 7$ spatial orbitals), non-orthogonal determinants built from
fragment-native orbitals~\cite{thom2009, sundstrom2014}, and general-$N$
non-orthogonal Slater--Condon evaluation by L\"owdin
cofactors~\cite{lowdin1955}.  Orbital exponents were optimized on
the \emph{isolated} Na atom only ($E(\mathrm{Na}) = -161.0845$~Ha against
a Clementi minimal-STO reference of $\approx -161.12$), with hydrogen at
$\zeta = 1$, so the molecular curve is a fragment prediction rather than
a fit.  Machinery validation at $R = 3.5\,a_0$: rotational invariance at
$1.1 \times 10^{-13}$, and the 91-determinant non-orthogonal complete space
reproduces the L\"owdin bitstring FCI to $1.8 \times 10^{-15}$~Ha.  (The
value serialized by the exploratory driver was corrupted by variable
shadowing---the record was written after a geometry sweep that rebinds the
nuclear-repulsion term---and has been corrected; the offset was exactly
$11/3.5 - 11/10 = 143/70$.)

\begin{table}[b]
\caption{NaH, all-electron 12e, $M = 7$.  Percentages are of in-basis FCI
binding.  Experiment: $R_{\rm eq} = 3.5667\,a_0$
($r_e = 1.8874$~\AA)~\cite{huberherzberg}.  The comparison value
$D_e = 1.961$~eV is inherited from the framework's earlier compilation and
its primary source is \emph{not yet verified}; it is not tabulated on the
NIST pages carrying $r_e$, so every $D_e$ percentage below should be read as
provisional pending that check.}
\label{tab:nah}
\begin{ruledtabular}
\begin{tabular}{lrrr}
method & $R_{\rm eq}$ ($a_0$) & $D_e$ (eV) & \% FCI \\
\hline
covalent (2 det) & 3.954 & 0.686 & 58.4 \\
$+\,\mathrm{Na^+H^-}$ (3 det) & \textbf{3.736} & \textbf{1.071} & \textbf{91.1} \\
all 4 det & 3.713 & 1.080 & 91.9 \\
FCI (91 det) & 3.595 & 1.175 & 100 \\
\hline
L\"owdin covalent (2 det) & --- & unbound & $-38.0$ \\
L\"owdin 3 det & 3.658 & 0.136 & 11.6 \\
\end{tabular}
\end{ruledtabular}
\end{table}

Three determinants give an interior minimum at $+4.8\%$ of the
experimental bond length, with the in-basis FCI landing at $+0.8\%$.  The
$D_e$ shortfall---$55\%$ of experiment retained, so a $45\%$ deficit---is
minimal-basis incompleteness (no polarization, no diffuse functions, and no
basis-set-superposition correction, the atom references being computed in
atom-only bases) and reproduces the pattern of the four-electron LiH case at
$53\%$; it is not an integral-fidelity effect.

As with LiH, the reading is sharper at \emph{fixed} geometry, where no
dissociation scan enters and the numbers are pinned by test.  At
$R = 3.5\,a_0$, against an isolated-atom dissociation reference of
$-161.584337$~Ha and an in-basis FCI binding of $1.170$~eV:\ two covalent
determinants recover $50.9\%$ of that binding, three recover $89.3\%$, and
all four $90.6\%$---while the \emph{same} three on L\"owdin-orthogonalized
orbitals recover $10.8\%$, and the L\"owdin two-determinant case is not
merely unbound but actively repulsive relative to the separated atoms
($-153\%$).  The well-minimum figures in Table~\ref{tab:nah} and these
fixed-geometry figures tell the same story by two routes.

The L\"owdin comparators are the informative rows.  Orthogonalizing the
orbitals and truncating to the \emph{same} number of determinants retains
$11.6\%$ of the binding, and the two-determinant case is destroyed
outright.  Because the complete space is bit-identically the same either
way, the damage is specific to \emph{truncated} spaces.  The four-electron
LiH case replicates this at $15.0\%$ and destroyed, respectively.

\begin{observation}[Orthogonalization damages only truncated spaces]
\label{obs:truncation}
Orthogonalization is span-preserving and therefore harmless on the
complete space, and severely harmful on truncated subspaces built from
physically meaningful fragment structures.  A qubit register only ever
represents a truncated space.
\end{observation}

Observation~\ref{obs:truncation} holds at four and twelve electrons,
first and second row, with and without $p$ orbitals---the four-electron
leg pinned by test, the twelve-electron leg on the exploratory driver.  We
expect it to generalize, and it is cheap to falsify.

The mechanism is also visible at \emph{fixed} geometry, without reference
to a dissociation limit, which makes it the cleanest form to check.  On
LiH at $R = 3.25\,a_0$ in a three-orbital fitted-Slater basis, against a
complete-space FCI of $-8.887049$~Ha: three fragment-native determinants
land $+1.32$~mHa above it, while the \emph{same three} determinants built
on L\"owdin-orthogonalized orbitals land $+40.52$~mHa above---a
$30.6\times$ larger error at identical determinant count.  At two
determinants the errors are $+19.4$~mHa and $+118.7$~mHa
($6.1\times$).  The complete spaces agree to machine precision, so the
damage is attributable to truncation and not to the orthogonalization
corrupting the integrals.

%% ========================================================================
\section{The quadrature-free diatomic}
\label{sec:qfd}
%% ========================================================================

\textbf{[MEASURED + SYMBOLIC]} The engine's classes have now been used
\emph{end-to-end}:\ complete diatomic FCI energies in which every one- and
two-electron integral is a closed-form expression---rationals, $\sqrt{\phantom.}$,
$\pi$, $\exp$, $E_1$, $\log$, $\gamma_E$ (the $\pi$'s are angular-measure and
normalization prefactors of the \emph{assembled} integrals;\ the radial seed content of
each class is the $\{\exp,E_1,\ln,\gamma_E\}$ set of the preceding sections, and
(AA$|$BB) is where even the prefactor $\pi$ cancels)---evaluated to arbitrary precision,
with no quadrature routine anywhere in the production path
(\texttt{geovac/qfd\_core.py}, \texttt{geovac/qfd\_assemble.py};\
\texttt{tests/test\_paper58\_qfd.py}).

Two pieces had to be added.  \emph{First}, the one-electron two-center closed
forms (overlap, kinetic, both nuclear-attraction kernels) did not previously
exist in the engine;\ they are derived from the Mulliken auxiliary integrals
$A_m(p)=\int_1^\infty \xi^m e^{-p\xi}d\xi$, $B_n(q)=\int_{-1}^{1}\eta^n
e^{-q\eta}d\eta$ in prolate spheroidal coordinates, where the Jacobian
factorization $(\xi^2-\eta^2)=(\xi+\eta)(\xi-\eta)$ renders both $1/r_A$ and
$1/r_B$ kernels polynomial;\ the kinetic energy uses the exact $s$-state radial
Laplacian, whose lowest surviving power is the $r^{-1}$ floor the auxiliary
form supports.  Two independent routes (explicit Laplacian vs the hydrogenic
eigen-trick) agree \emph{symbolically---difference exactly~$0$}---on
homonuclear and heteronuclear sets.  \emph{Second}, the exchange assembly is
re-run keeping every factor symbolic (the library entry point casts to float,
capping certification at $\sim\!16$ digits).

\textbf{[MEASURED]} H$_2$ ($1s$ per center, $\zeta=1$, $R=1.4$):\ the
homonuclear exchange Neumann $\tau$-series \emph{terminates symbolically} \textbf{[SYMBOLIC]}
($\tau=1$ and all $\tau>2$ are exact zeros), so the pipeline contains no
truncation of any kind, and the FCI total energy is certified to \textbf{60
digits} (under the reference artifact's capping convention---a claim never exceeds the
weaker run's requested precision;\ the dps~$60$ vs $90$ evaluations in fact agree to
$6\times10^{-85}$, so the recorded value is conservative):
\[
E(1.4,\ \zeta{=}1)=-1.10655660609135850801940122475993772281832890689690719549979\ldots
\]
with $R=1.6$, $2.0$ and the variational $\zeta=1.197$ point certified
alongside.  Every integral matches an \emph{independent} literature closed
form to $\sim\!10^{-61}$---the exchange integral against
Sugiura~\cite{sugiura1927}, whose expression independently carries $\ln$,
Euler's $\gamma$ and $\mathrm{Ei}$, corroborating this paper's
$\{E_1,\ln,\gamma\}$ seed set for that class from a route nearly a century
removed---and matches direct quadrature at $10^{-21}$--$10^{-23}$.

\textbf{[MEASURED]} LiH ($s$-only minimal:\ Li~$1s,2s$ + H~$1s$, $R=3.015$):\
the heteronuclear $\tau$-series does \emph{not} terminate, so the assembly is
zero-quadrature but not truncation-free;\ the certified count is taken
\emph{net of an explicit tail bound} (measured monotone $\tau$-ratios give
$|{\rm tail}|\le 1.95\times10^{-32}$ summed over quartets, propagated through
the 2-RDM weights and $\|S^{-1/2}\|^4$ to $1.2\times10^{-30}$~Ha):\
$E_{\rm tot}=-7.87261979241561217317558085835\ldots$~Ha,
\textbf{30 digits}.  Honest scope, per the framework's own doctrine:\ these
are \emph{minimal-basis} numbers---the certified $\zeta{=}1$ H$_2$ value sits
$0.068$~Ha above exact ($0.027$~Ha at the variational $\zeta=1.197$ point) and
the $s$-only LiH far above $-8.0705$---and what is certified is the assembly
and its precision, not accuracy;\ the accuracy axis remains basis size.
Certified values, with per-integral transcendence tags, are collected in the
reference artifact (\texttt{docs/certified\_reference\_values.md},
\texttt{benchmarks/certified\_reference/}).

\textbf{[SYMBOLIC + MEASURED]} The closed forms compose one level further:\ the
\emph{entire minimal-basis H$_2$ potential curve is a single explicit
expression}.  In the $1s/1s$ basis the singlet-$\Sigma_g^+$ CI space is
$\{\,|g^2\rangle,\,|u^2\rangle\,\}$, so the FCI ground state is the
quadratic-formula root of a $2\times2$ matrix whose entries are the closed-form
molecular-orbital integrals,
\begin{equation}
E(R)=\tfrac12\!\left(H_{11}+H_{22}\right)
-\sqrt{\tfrac14\!\left(H_{11}-H_{22}\right)^{2}+K_{gu}^{2}}+\frac1R,
\label{eq:h2_pes}
\end{equation}
with $H_{11}=2h_{gg}+J_{gg}$, $H_{22}=2h_{uu}+J_{uu}$, and every ingredient an
explicit closed form in $R$---the assembled expression contains exactly the
atoms $\{\exp,\,E_{1},\,\ln,\,\gamma_{E}\}$ and nothing else beyond the rational and normalization prefactors of the opening list.  A potential
curve as a formula, not a point table:\ its exact $R$-derivatives certify the
equilibrium constants by closed-form Newton iteration
(\texttt{geovac/qfd\_assemble.py::h2\_closed\_form\_E};\
\texttt{tests/test\_paper58\_qfd.py}),
\begin{align*}
R_{\rm eq}&=1.66799996697274872492704641030726451334980414\ldots\ \text{bohr},\\
D_e&=0.118650362098295311853497764333884671691909683\ldots\ \text{Ha},\\
k=E''(R_{\rm eq})&=0.254703934307969981152727935724355672678251789\ldots\ \text{Ha/bohr}^2,
\end{align*}
(Newton at dps $45$ vs $60$ agreeing to $9.5\times10^{-46}$;\ the curve
dissociates to exactly $-1$~Ha---two $\zeta{=}1$ hydrogen atoms---verified at
$3\times10^{-49}$, so $D_e=-1-E(R_{\rm eq})$ is itself closed-form).  Honest
frame:\ these are the minimal-basis constants ($R_{\rm eq}$ exact is $1.401$,
$k$ maps to $\omega_e\approx3.7\times10^{3}$~cm$^{-1}$ vs the measured
$4401$---the conversion importing the reduced mass is Layer-2;\ and the
minimal-basis H$_2$ CI itself is as old as Sugiura's integral);\ what is new is
their \emph{status}:\ equilibrium geometry and force constant of a two-electron
molecule as certified roots of one closed-form expression.

Two engine findings from the build.  (i)~\emph{Domain restriction}
(measured):\ the shell-route hybrid evaluator with $l_a,l_b>0$ requires
$Z_B<Z_A$ strictly---at $Z_B=Z_A$ a removable rate coincidence produces
NaN (the limit exists and matches quadrature), and $Z_B>Z_A$ reaches an
unimplemented $\mathrm{Ei}$ branch;\ a coverage gap, not a wrong value.
(ii)~A pre-existing exploratory driver's ``DISAGREE beyond the fit
floor---investigate'' verdict against a Gaussian reference is resolved
\emph{against the reference}:\ the $8\times10^{-6}$ gap sits in the Gaussian
one-electron block (the cusp), while the native block matches literature
closed forms to $60$ digits.

\section{Scope}
\label{sec:scope}
%% ========================================================================

What this paper does not claim.  It does not claim the framework's
atomic-sector results are affected: single-center angular structure is
untouched, exact, and remains the basis of the atomic qubit results.  It
does not claim the qubit-resource results of Paper~14 are wrong; it
clarifies that they are measured on the builder, whose neglected overlap
is the subject here.  And it does not claim a route around the
obstruction: the census closes the one route the roadmap had left---%
untransformed bare-tensor use---while
Observation~\ref{obs:cost_conservation} records only that the five
relocations attempted so far each moved the cost without reducing it.

A distinction worth naming explicitly, because it bounds any scaling
ambition.  \emph{Symmetry sparsity}---exact zeros forced by conservation
laws, knowable from labels---is what this framework has, and
Theorem~\ref{thm:abelian_residue} shows it is an atomic-sector property
that degrades to an abelian residue at two centers and to a finite
grading when linearity is lost.  \emph{Distance sparsity}---approximate
near-zeros from spatial locality, exploited by screening, density
fitting, and density-matrix truncation---is what linear-scaling
electronic structure runs on, and it \emph{improves} with system size.
The framework has the first and none of the second.  It is, structurally,
a small-system framework, and both of its advantages contract as atom
count grows: the abelian grading shrinks with symmetry, while the dense
cross-center classes grow with the number of center pairs.

\begin{table}[b]
\caption{Claim $\to$ backing.  \textsc{live} $=$ passing test in the
\texttt{tests/test\_paper58\_*} suite (six files, plus
\texttt{tests/test\_two\_center\_eri\_aabb.py} for the $(AA|BB)$ rows; per-claim mapping in
\texttt{docs/claim\_test\_matrix.md}).  \textsc{owed} $=$
result stands on an exploratory driver and must be migrated
self-contained into \texttt{tests/} before this paper is a finished
artifact.  \textsc{corrob.} $=$ numerically corroborated but not exactly
decided; see Sec.~\ref{sec:scope}.}
\label{tab:backing}
\begin{ruledtabular}
\begin{tabular}{lll}
claim & tier & backing \\
\hline
Table~\ref{tab:census} $S$, $h$ & MEASURED & \textsc{live} \\
Table~\ref{tab:census} $g$ & COUNTED & \textsc{corrob.} \\
2-center ERI engine, per class & MEASURED & \textsc{live} \\
axial invariance & MEASURED & \textsc{live} \\
Eq.~\eqref{eq:braket_certificate} & SYMBOLIC & \textsc{live} \\
Thm.~\ref{thm:abelian_residue}(i) & INT. & \textsc{live} \\
Thm.~\ref{thm:abelian_residue}(ii) & INT. & \textsc{live} \\
2-center machinery & SYMB. & \textsc{live} \\
Obs.~\ref{obs:no_angle} & MEASURED & \textsc{live} \\
Pred.~\ref{pred:angular} & UNTESTABLE & (needs angular build) \\
Obs.~\ref{obs:cost_conservation} & OBSERVATION & CHANGELOG v4.73.x \\
Obs.~\ref{obs:permutation_cost} & MEASURED & \textsc{live} \\
Tab.~\ref{tab:nah}, fixed $R$ & MEAS. & \textsc{live} \\
Tab.~\ref{tab:nah}, minima & PANEL & \textsc{owed} \\
NaH atom-only $\zeta$ & MEAS. & \textsc{live} \\
Obs.~\ref{obs:truncation} & MEASURED & \textsc{live} \\
fit quality, span id. & MEASURED & \textsc{live} \\
QFD 1e closed forms (Sec.~\ref{sec:qfd}) & SYMBOLIC & \textsc{live} \\
QFD H$_2$ 60-digit energy & MEASURED & \textsc{live} (fast legs) + certified tables \\
QFD LiH 30-digit ($\tau$-tail net) & MEASURED & \textsc{owed} (campaign value; fast legs pin one exchange quartet only) \\
\end{tabular}
\end{ruledtabular}
\end{table}

%% ========================================================================
\begin{thebibliography}{99}

\bibitem{sugiura1927}
Y.~Sugiura, ``\"Uber die Eigenschaften des Wasserstoffmolek\"uls im
Grundzustande,'' Z.~Phys.\ \textbf{45}, 484--492 (1927).

\bibitem{loutey_paper11}
J.~Loutey, ``The Molecular Fock Projection:\ Diatomic Molecules as
Prolate Spheroidal Lattices,'' GeoVac Paper~11 (2026).

\bibitem{loutey_paper14}
J.~Loutey, ``Structurally Sparse Qubit Hamiltonians from Spectral Graph
Theory,'' GeoVac Paper~14 (2026).

\bibitem{loutey_paper17}
J.~Loutey, ``Composed Natural Geometries for Core-Valence Diatomics:\
LiH and BeH$^{+}$ from Discrete Graph Topology,'' GeoVac Paper~17 (2026).

\bibitem{loutey_paper19}
J.~Loutey, ``Toward PK-Free Molecular Hamiltonians:\ Two-Center
Sturmian Integrals and Coupled Composition,'' GeoVac Paper~19 (2026).

\bibitem{loutey_paper20}
J.~Loutey, ``Structurally Sparse Molecular Hamiltonians for
Near-Term Quantum Simulation,'' GeoVac Paper~20 (2026).

\bibitem{huberherzberg}
K.~P.~Huber and G.~Herzberg, \textit{Constants of Diatomic Molecules}
(Van Nostrand Reinhold, New York, 1979);\ $^{23}$NaH $X\,^1\Sigma^+$
$r_e = 1.8874$~\AA\ $= 3.5667\,a_0$, as tabulated in the NIST Chemistry
WebBook (accessed 2026-08-11).

\bibitem{mulliken1949}
R.~S.~Mulliken, C.~A.~Rieke, D.~Orloff, and H.~Orloff, ``Formulas and
numerical tables for overlap integrals,''
\textit{J.\ Chem.\ Phys.}\ \textbf{17}, 1248--1267 (1949).

\bibitem{roothaan1951}
C.~C.~J.~Roothaan, ``A study of two-center integrals useful in
calculations on molecular structure.\ I,''
\textit{J.\ Chem.\ Phys.}\ \textbf{19}, 1445--1458 (1951).

\bibitem{ruedenberg1951}
K.~Ruedenberg, ``A study of two-center integrals useful in calculations
on molecular structure.\ II.\ The two-center exchange integrals,''
\textit{J.\ Chem.\ Phys.}\ \textbf{19}, 1459--1477 (1951).

\bibitem{lowdin1950}
P.-O.~L\"owdin, ``On the non-orthogonality problem connected with the use
of atomic wave functions in the theory of molecules and crystals,''
\textit{J.\ Chem.\ Phys.}\ \textbf{18}, 365--375 (1950).

\bibitem{mcmurchie1978}
L.~E.~McMurchie and E.~R.~Davidson, ``One- and two-electron integrals
over Cartesian Gaussian functions,''
\textit{J.\ Comput.\ Phys.}\ \textbf{26}, 218--231 (1978).

\bibitem{artacho1991}
E.~Artacho and L.~Mil\'ans del~Bosch, ``Nonorthogonal basis sets in
quantum mechanics:\ Representations and second quantization,''
\textit{Phys.\ Rev.\ A}\ \textbf{43}, 5770--5777 (1991).

\bibitem{avery1989}
J.~Avery, \textit{Hyperspherical Harmonics:\ Applications in Quantum
Theory} (Kluwer Academic, Dordrecht, 1989).

\bibitem{avery2006}
J.~E.~Avery and J.~S.~Avery, \textit{Generalized Sturmians and Atomic
Spectra} (World Scientific, Singapore, 2006).

\bibitem{shibuya1965}
T.~Shibuya and C.~E.~Wulfman, ``Molecular orbitals in momentum space,''
\textit{Proc.\ R.\ Soc.\ Lond.\ A}\ \textbf{286}, 376 (1965).

\bibitem{lowdin1955}
P.-O.~L\"owdin, ``Quantum theory of many-particle systems.\ I.\ Physical
interpretations by means of density matrices, natural spin-orbitals, and
convergence problems in the method of configurational interaction,''
\textit{Phys.\ Rev.}\ \textbf{97}, 1474--1489 (1955).

\bibitem{thom2009}
A.~J.~W.~Thom and M.~Head-Gordon, ``Hartree--Fock solutions as a
quasidiabatic basis for nonorthogonal configuration interaction,''
\textit{J.\ Chem.\ Phys.}\ \textbf{131}, 124113 (2009).

\bibitem{sundstrom2014}
E.~J.~Sundstrom and M.~Head-Gordon, ``Non-orthogonal configuration
interaction for the calculation of multielectron excited states,''
\textit{J.\ Chem.\ Phys.}\ \textbf{140}, 114103 (2014).

\bibitem{barnett_coulson1951}
M.~P.~Barnett and C.~A.~Coulson, ``The evaluation of integrals occurring
in the theory of molecular structure.\ Parts I \& II,''
\textit{Philos.\ Trans.\ R.\ Soc.\ Lond.\ A}\ \textbf{243}, 221--249 (1951).

\bibitem{filter_steinborn1978}
E.~Filter and E.~O.~Steinborn, ``Extremely compact formulas for molecular
two-center one-electron integrals and Coulomb integrals over Slater-type
atomic orbitals,'' \textit{Phys.\ Rev.\ A}\ \textbf{18}, 1 (1978).

\end{thebibliography}

\end{document}
